\chapter{是谁证明 DNA 携带遗传信息}

\begin{kaichang}
如果你或家人做过基因检测，多半见过那样一支小管子：吐一口唾液进去，寄走，几周后收到一份报告，告诉你祖源构成、乳糖耐受能力、酒精代谢快慢。亲子鉴定也差不多——用棉签在口腔内壁刮几下，送去实验室，就能判断两个人是不是亲生父子。

先停下来想一想这里面的一个巨大跳跃：凭什么？凭什么一口唾液里藏着的某种分子，就能“代表你”？你的血型、卷不卷舌头、容不容易喝醉，凭什么由它来记账？

今天你觉得答案天经地义：DNA 嘛，谁都知道。可在 20 世纪上半叶，这个答案不但不明显，而且几乎没人愿意相信。当时全世界最顶尖的生物化学家，绝大多数把宝押在另一种分子身上——蛋白质。这一章就来讲：科学界怎样一个实验一个实验地，把答案从“多半是蛋白质”掰回到“就是 DNA”。你会看到，没有任何一个实验能一锤定音，真正把结论钉死的，是一条互相咬合的证据链。
\end{kaichang}

\section{一场几乎所有人都押错宝的赌局}

先把镜头拉回 1940 年代的实验室，只描述当时测得到的东西。

生物学家早就知道，细胞核里有一种线状结构叫染色体：细胞分裂前它会复制一份，再精确地平分给两个子细胞（第 64 章讲过这套搬运流程）。第 63 章孟德尔的豌豆实验又告诉大家：遗传信息是分成一个个“因子”传递的。把两件事拼在一起，结论很自然——遗传信息就装在染色体上。

那么染色体是用什么做的？化学分析给出两种主要成分。一种是蛋白质：第 48 章认识过它，由 20 种氨基酸连成，能折叠成千变万化的形状，当时已知的酶、抗体、许多激素全是蛋白质。另一种叫脱氧核糖核酸，简称 DNA：由核苷酸连成，而核苷酸的差别只在 4 种含氮碱基——腺嘌呤（A）、鸟嘌呤（G）、胞嘧啶（C）、胸腺嘧啶（T）。就 4 种。

现在请你站在当年的立场上下注：谁是遗传信息的载体？绝大多数科学家押蛋白质，理由听起来很充分。信息意味着千变万化，蛋白质有 20 种“字母”，排列组合丰富无比；DNA 只有 4 种字母，而且当时流行的“四核苷酸假说”认为，DNA 不过是 4 种核苷酸等比例地一遍遍重复，单调乏味，最多给染色体当个骨架。这个假说来自化学家列文早年的粗略测量——他测出的碱基比例恰好接近一比一，于是“单调重复”的图景就这样以讹传讹地写进了教科书，一写就是几十年。20 种字母对 4 种字母，还用比吗？

这个推理错在哪里，本章中段会揭晓。先记住一点：科学史上“大家都这么认为”，从来不等于“证据支持这样认为”。当时并没有任何一个实验直接比较过两种分子的遗传能力——人们只是凭“谁看起来更复杂”投了票。

\section{死细菌改变了活细菌的“品种”}

第一道裂缝出现在 1928 年的伦敦。当时距一战结束不久，肺炎仍是致命的常见病，英国卫生部想知道不同菌株的致病能力为何不同，好为疫苗做准备——格里菲思（Frederick Griffith）正是在这项任务里研究肺炎链球菌，一种能引起肺炎的细菌。这种细菌有两个品种，肉眼就能区分：S 型（光滑型）菌落光滑发亮，外面裹着一层多糖荚膜，能躲过小鼠的免疫系统，会致病；R 型（粗糙型）没有荚膜，菌落粗糙，不致病。

格里菲思做了四组注射实验：

\begin{table}[htbp]
\centering
\begin{tabular}{lll}
\toprule
注射进小鼠的内容 & 小鼠结局 & 死后分离到的活菌 \\
\midrule
活 R 型 & 存活 & 无 \\
活 S 型 & 死亡 & S 型 \\
加热杀死的 S 型 & 存活 & 无 \\
活 R 型 + 加热杀死的 S 型 & 死亡 & 活的 S 型 \\
\bottomrule
\end{tabular}
\end{table}

前三组都不意外。第四组才真正炸开了锅：单独注射活 R 型，小鼠活；单独注射死 S 型，小鼠也活；可两者混在一起注射，小鼠死了，而且从尸体里分离出了活的 S 型细菌。这些 S 型还能继续繁殖，后代统统是 S 型。

先澄清一个容易想岔的地方：死细菌并没有“复活”——加热杀死的 S 型单独注射时，小鼠活得好好的。真正发生的事情是：少数活 R 型细菌吸收了死 S 型释放出来的某种东西，获得了制造荚膜的“配方”，从此自己和后代都变成 S 型。复活的不是细菌个体，而是一段信息。

格里菲思把这种现象叫“转化”，把那个看不见的东西叫“转化因子”。现在来做本章反复要做的训练：这个实验\textbf{能证明什么}？它证明了细菌体内存在一种可以转移、可以遗传的物质，能从分子层面永久改变一个生物的性状——“遗传物质”第一次被实验逼到台前。但它\textbf{不能证明}转化因子是哪种分子。格里菲思本人倾向于认为那只是某种蛋白质或荚膜物质。他把门推开了一条缝，自己却没有走进去。

\section{把候选分子一个个“请出考场”}

真正把接力棒捡起来的是洛克菲勒研究所的艾弗里（Oswald Avery）和同事麦克劳德、麦卡蒂。他们花了十几年，先在试管里重现转化（不再依赖小鼠），然后做了一件逻辑上极其漂亮的事：\textbf{排除法}。

他们把 S 型细菌打碎、提纯出内含物，再逐一“请出”候选分子，看转化能力是否还在：

\begin{itemize}[itemsep=2pt]
\item 用蛋白酶处理，把蛋白质降解掉——转化照样发生；
\item 用 RNA 酶处理，把 RNA 降解掉——转化照样发生；
\item 用能分解荚膜多糖的酶处理——转化照样发生；
\item 用 DNA 酶处理，专门把 DNA 切断——转化彻底消失。
\end{itemize}

1944 年，他们发表结论：转化因子是 DNA，不只是结构成分，而是真正起作用的遗传物质。

这个实验能证明什么？在高度纯化的条件下，DNA 本身就足以把遗传性状从一个细菌转移给另一个细菌——蛋白质、RNA、多糖都不在场时，信息依然在起作用。逻辑无懈可击吗？几乎。当时有批评者咬住一点：提取物再纯，谁敢保证没有百万分之几的蛋白质残留？也许真正起作用的是那一点没除干净的蛋白质。这个质疑今天看像钻牛角尖，但在“人人都押蛋白质”的年代，它足以让多数人继续观望。艾弗里本人又极其谨慎，论文措辞克制，结果这项后来被认为值得诺贝尔奖的工作，当时只换来礼貌的沉默。你看，证据的分量不仅取决于实验本身，还取决于时代准备好没有。

\section{给两种分子贴上不同的“发光标签”}

真正把门撞开的是 1952 年的赫尔希（Alfred Hershey）和蔡斯（Martha Chase）。他们换了一个更干净的系统：T2 噬菌体，一种专门感染细菌的病毒。说它“干净”，是因为它身上只剩两种成分——一个蛋白质外壳，里面装着一条 DNA，候选分子从一整只细菌的成千上万种分子骤减到两种。它感染细菌的方式像打针：外壳贴在细菌表面，把“内容物”注入细菌，细菌工厂随即被迫批量生产新噬菌体，最后细菌破裂，放出成百上千个子代。

关键问题来了：被注进细菌的到底是 DNA 还是蛋白质？谁进去了，谁就是遗传信息——因为子代噬菌体的外壳和 DNA，都是按“进入者”的指令造出来的。

赫尔希和蔡斯的妙招，是给两种分子贴上不同的放射性标签。第 4 章讲过同位素：同一元素的不同同位素化学行为完全一样，但放射性同位素会发出可以探测的射线。而两种分子的元素组成恰好错开——蛋白质含硫却几乎不含磷（有两种氨基酸含硫），DNA 的糖—磷酸骨架含大量磷却完全不含硫。于是：

\begin{itemize}[itemsep=2pt]
\item 让噬菌体在含放射性硫 $^{35}\mathrm{S}$ 的环境里繁殖——只有蛋白质外壳被标记；
\item 让噬菌体在含放射性磷 $^{32}\mathrm{P}$ 的环境里繁殖——只有 DNA 被标记。
\end{itemize}

然后让标记过的噬菌体去感染细菌，几分钟后用一台普通的厨房搅拌机猛搅，把吸附在细菌表面的空外壳甩掉，再离心：细菌沉到管底，轻的外壳留在上清液里。测放射性在哪一层，就知道谁进了细胞：

\begin{table}[htbp]
\centering
\begin{tabular}{llll}
\toprule
标记 & 被标记的分子 & 放射性主要在哪 & 子代噬菌体 \\
\midrule
$^{35}\mathrm{S}$ & 蛋白质外壳 & 上清液（外壳） & 检测不到放射性 \\
$^{32}\mathrm{P}$ & DNA & 细菌沉淀中 & 部分子代带 $^{32}\mathrm{P}$ \\
\bottomrule
\end{tabular}
\end{table}

结论一目了然：进入细菌的是 DNA，蛋白质外壳被留在门外当了“空针筒”；而进入的 DNA 指导细菌造出了成百上千个完整的子代噬菌体——包括它们全新的蛋白质外壳。发号施令的是 DNA。

再来做一次训练：这个实验\textbf{单独}能证明一切吗？不能。仍有少量 $^{35}\mathrm{S}$ 随细菌沉淀，批评者可以说“也许有微量蛋白质溜进去了”；而且它只研究了噬菌体这一种生物。它真正不可替代的贡献，是在一个“进出分明”的系统里正面追踪了两种分子各自的行踪。艾弗里的软肋是“提取物里可能残留蛋白质”，赫尔希和蔡斯恰好把它变成“进门的到底是谁”的直接目击——两个实验的漏洞互相堵上。这就是为什么教科书把它们并排写：证据不靠一个英雄实验，而靠链条咬合。

\begin{wuqu}
“赫尔希和蔡斯的搅拌机实验一锤定音，单独证明了 DNA 是遗传物质。”——这是流传最广的简化。单独看，它既不能排除少量蛋白质进入细菌，也只覆盖噬菌体；真正让人无法反驳的，是它与艾弗里的酶解实验互相补位：一个证明“纯 DNA 足以传递性状”，一个证明“传递性状时进门的是 DNA”。科学里“一锤定音”的故事大多是事后改写出来的，真实的过程是多个实验从不同角度逼近同一个结论，直到替代解释被一个个挤出去。
\end{wuqu}

\section{四种字母的对账本}

噬菌体实验回答了“谁进门”，但还有一笔旧账没了：DNA 真的那么单调吗？如果它永远 4 种碱基等比例重复，那它确实存不了多少信息——四核苷酸假说不拆，蛋白质派就还有退路。

拆这个台的是生物化学家查加夫（Erwin Chargaff）。1940 年代末，他用纸层析技术把不同生物的 DNA 拆开，精确称量 4 种碱基各占多少。结果有两条。

第一条：不同物种的 DNA，碱基比例各不相同。人的、牛的、细菌的、酵母的，都不一样。DNA 不是千篇一律的单调骨架，它的“配方”随物种而变——这正是携带信息的前提。

第二条更奇怪：无论测什么物种，A 的含量总约等于 T，G 总约等于 C。这就是著名的“查加夫规则”。例如他测人的 DNA：A 约占 30.9\%，T 约占 29.4\%，G 约占 19.9\%，C 约占 19.8\%——并不分毫不差（测量总有误差），但 A 与 T、G 与 C 的接近绝不是巧合。查加夫本人没参透这个等式的含义，可它等于在大声宣布：DNA 内部的碱基是\textbf{成对}的。1952 年沃森和克里克登门拜访时，查加夫当面讲过这些数据；几年后，这个等式成了拼出双螺旋的关键拼图片。

\begin{qiaoliang}
回到本章开头那场赌局：4 种字母真的太少吗？算一笔账。一段序列有 $n$ 个位置，每个位置可以是 A、T、G、C 之一，总共就有 $4^n$ 种不同序列。$n=10$ 时约 100 万种；$n=100$ 时约 $10^{60}$ 种；一段 3000 个碱基的序列有 $4^{3000}\approx 10^{1800}$ 种写法——而整个可观测宇宙的原子总数也不过约 $10^{80}$ 个。4 种字母的信息容量大到荒谬。蛋白质派当年的直觉错不在数学，而在于被四核苷酸假说误导，以为 DNA 的 4 种字母只能等比例重复、排列没有自由。查加夫的数据恰恰证明：比例随物种而变，排列是自由的——4 种字母完全够写一部遗传之书。
\end{qiaoliang}

\section{照片 51 号：功劳、竞争与底线}

\begin{zhengju}
最后一块拼图是 DNA 的形状。1950 年代初，好几个实验室正在竞赛破解它的三维结构。伦敦国王学院的罗莎琳德·富兰克林（Rosalind Franklin）是当时最出色的 X 射线衍射专家之一。X 射线衍射的原理可以这样理解：让 X 射线穿过排列整齐的分子样品，射线被规则排列的原子散射，形成特定图样；从斑点的位置可以反推分子的几何形状——有点像从栅栏影子的花纹，反推栅栏本身的间距。

1952 年，富兰克林和研究生戈斯林拍下了著名的“照片 51 号”：一张 DNA 纤维的衍射图，中央一个清晰的 X 形。懂行的人一眼就能看出，X 形是螺旋结构的衍射签名，图样里还藏着螺旋的直径和每一圈的升高距离。富兰克林测出的数据还表明，DNA 的亲水骨架在分子的外侧。

1953 年 1 月，她的同事威尔金斯在没有征得她同意的情况下，把照片 51 号拿给来访的沃森看。沃森后来写道，看到照片的瞬间“我的嘴巴张得老大”——螺旋的猜测被证实了。与此同时，富兰克林整理在一份内部报告里的测量数据，也辗转到了克里克手中。几周之内，沃森和克里克拼出了双螺旋模型：两条链反向缠绕，A 与 T 配对、G 与 C 配对——正好解释查加夫的等式。1953 年 4 月，论文发表，同期也刊登了富兰克林和威尔金斯各自的数据论文。

接下来是这个故事必须认真讲的部分。富兰克林当时已离开国王学院，转而研究烟草花叶病毒，成果同样出色；1958 年，她因卵巢癌去世，年仅 37 岁。1962 年，诺贝尔生理学或医学奖授予沃森、克里克和威尔金斯——诺贝尔奖只授予在世者，她连被提名的机会都没有。而沃森在回忆录《双螺旋》里，把她描绘成一个难相处、不懂自己数据价值的配角，许多读者直到多年后读了整理的史料才知道：照片 51 号是她拍的，关键数据是她测的。

该怎样评价这段历史？科学史不做“如果”题——我们无法假设富兰克林再多一年会不会自己解出结构。但底线是清楚的：\textbf{数据的使用应当征得同意，贡献应当被如实记录}。合作和竞争都能推动科学，可“谁先看到照片”不该遮蔽“谁拍出了照片”。今天，几乎所有教科书都会把富兰克林的名字与双螺旋写在一起，这不是施舍，而是科学共同体在补记一笔迟到的账。
\end{zhengju}

\section{一张表看懂整条证据链}

走到这里，该退一步，把五个环节摆在一起看清楚——每一步回答了什么、又留下了什么：

\begin{table}[htbp]
\centering
\small
\begin{tabular}{p{3.0cm}p{4.8cm}p{4.8cm}}
\toprule
实验或发现 & 它能证明 & 它不能单独证明 \\
\midrule
格里菲思（1928） & 存在可转移、可遗传的“转化因子” & 转化因子是哪种分子 \\
艾弗里（1944） & 高度纯化的 DNA 足以传递性状 & 提取物绝无微量杂质；对其他生物也成立 \\
赫尔希—蔡斯（1952） & 噬菌体感染时进入细菌的是 DNA & 完全没有蛋白质混进去；所有生物都用 DNA \\
查加夫（1949 起） & DNA 碱基比例因物种而异，且 A$\approx$T、G$\approx$C & DNA 的具体三维结构 \\
富兰克林与双螺旋（1953） & DNA 的螺旋几何与碱基配对方式 & 结构怎样指导复制（下一章讲） \\
\bottomrule
\end{tabular}
\end{table}

这张表就是本章的边界层：每个实验单独看都有缝隙，整条链子合起来才密不透风。它也是本书第五个贯穿问题——“在生命中，信息怎样被保存、复制、调节和选择”——的第一块基石：在回答“怎样保存、怎样复制”之前，必须先确认“保存在什么分子上”。这个确认来之不易，花了整整四分之一个世纪。还要再画一条边界：说“DNA 是遗传物质”，准确的说法是“DNA 是\textbf{主要的}遗传物质”。1956 年，科学家把烟草花叶病毒拆开又重组，证明这种病毒的遗传物质是 RNA，根本不需要 DNA；你今天听说的流感病毒、新冠病毒，也都是 RNA 病毒。另外，你可能听说过“疯牛病的朊病毒靠蛋白质遗传”——朊病毒确实是一种会“传染”的蛋白质，但它不携带序列信息，只是诱导正常蛋白质折叠成同样的错误形状，并不能像 DNA 或 RNA 那样记录千变万化的信息。把它叫“遗传物质”，是对证据的过度引申。

\begin{shiyan}
DNA 听起来高深，其实你在家就能把它从水果里“拉”出来。材料：熟透的香蕉半根（或草莓两颗）、食盐一小勺、洗洁精、温水、纱布或咖啡滤纸、透明杯子、冰镇的医用酒精（或高度白酒）。

步骤：① 香蕉装进保鲜袋压成泥，加盐水（一小勺盐溶于半杯温水）和几滴洗洁精，轻轻搅匀，静置 10 分钟——洗洁精负责溶解细胞膜和核膜的磷脂双层（第 51 章讲过），盐帮助 DNA 聚集；② 用纱布过滤，留滤液；③ 将冰酒精沿杯壁缓缓倒入，让它浮在滤液上，不要搅动；④ 静置几分钟，观察两层交界处。

观察点：交界处会慢慢出现白色絮状物，像云雾一样，可以用牙签绕出来——那就是成团的 DNA。注意，你看到的是几万亿条 DNA 分子缠成的絮团；单个双螺旋只有 2 纳米宽，任何光学显微镜都看不见它。

安全提示：酒精易燃，全程远离明火；接触过洗洁精的材料不可再吃；洗洁精和酒精都不要弄进眼睛。
\end{shiyan}

\section{回到那管唾液}

现在可以回答开场的问题了。你吐进采样管的唾液里，混着从口腔内壁脱落的细胞；每个细胞的细胞核里都有 46 条染色体，上面盘着 DNA。你的这份 DNA 一半来自父亲、一半来自母亲——减数分裂怎样完成这次“对半分”，第 64 章讲过。几十亿个碱基对里，大约有千分之一的字母与别人不同，这些差异就是亲子鉴定、刑侦比对、基因检测读取的全部内容。实验室做的事，本质上就是读出这些字母，再比一比谁和谁更像。

感受一下这本“书”的篇幅。一个普通人体细胞里的 DNA 拉直约 2 米，含约 60 亿个碱基对；你全身约有 $3\times 10^{13}$ 个细胞，所有 DNA 首尾相接，总长 $6\times 10^{13}$ 米，也就是 600 亿千米。地球到太阳的距离约 $1.5\times 10^{11}$ 米，除一除：相当于 400 个日地距离，足够在地球和太阳之间往返 200 趟。而这一切，都是从格里菲思那只小鼠开始，被一节一节实验链条确认为“信息载体”的东西。

这条证据链也早已走出教科书：DNA 比对帮助侦破和昭雪了无数案件——刑侦实验室并不读完整个基因组，只比对十几到二十个因人而异的“短串联重复”位置，两个人全部吻合的概率就低到几乎不可能；核酸检测在疫情期间成了亿万人熟悉的词，它检测的其实是病毒自己的核酸序列，从另一条路线再次证明“核酸携带信息”；而读懂和改写 DNA 的技术，正在变成医学和农业的日常工具——第 82 章会专门讲“改写基因”这件事能做到什么、又该在哪里停下。

\section{信息找到了，读法还没找到}

本章结束时，“谁是遗传物质”已经尘埃落定，但一个更深的问题刚刚露头：DNA 的结构凭什么既能\textbf{稳定保存}几十亿个字母，又能在细胞分裂前\textbf{精确复制}一份？要知道，同样是写在纸上的信息，印刷术和手抄本的出错率天差地别——分子世界的“抄写”凭什么做到几十亿个字母几乎不出错？查加夫的 A 等于 T 等式，和照片 51 号里那个 X 形，已经把答案暴露了一半：A 和 T 为什么正好成对？两条链为什么反向缠绕？复制时，碱基配对怎样自动保证抄写的准确？下一章《DNA 的结构为什么适合保存信息》，我们把双螺旋拆开，看它怎样同时扮演“保险柜”和“复印机”。

\begin{tiaozhan}
赫尔希—蔡斯的标签为什么能做到“只标蛋白质、只标 DNA”？关键在元素分布：20 种氨基酸中只有甲硫氨酸和半胱氨酸含硫，而 DNA 完全不含硫；DNA 的糖—磷酸骨架上每个核苷酸都带一个磷原子，蛋白质却几乎不含磷。所以用 $^{35}\mathrm{S}$ 培养，标记只会落在蛋白质外壳上；用 $^{32}\mathrm{P}$ 培养，标记只会落在 DNA 上。探测上，$^{32}\mathrm{P}$ 的半衰期约 14.3 天，$^{35}\mathrm{S}$ 约 87 天，衰变放出的射线能让计数器或感光胶片产生信号——只要实验在几周内完成，标签就始终“亮着”。还有一笔值得一提的史事：赫尔希因这项工作分享了 1969 年诺贝尔奖，而实验的主要动手者蔡斯，当年仅以技术员身份署名。科学史里的署名问题，不止富兰克林一例。
\end{tiaozhan}

\section*{练一练}

\begin{sikaoti}
\item 画一幅信息流图：格里菲思第四组实验中，“制造荚膜的信息”从哪里出发、经过什么载体、最终在哪里表达？用方框和箭头表示，并注明：信息没有随死细菌消失。
\item 假如你是 1944 年质疑艾弗里的科学家，认定“残留蛋白质才是真凶”，你会设计一个什么新实验来检验自己的质疑？（提示：赫尔希和蔡斯八年后是怎样做的？）
\item 赫尔希和蔡斯为什么不能用碳或氮的标记来区分蛋白质和 DNA？（提示：想想两种分子各自含有哪些元素。）
\item 查加夫测得某种生物 DNA 中 A 占 24\%，那么 T 约占多少？G 和 C 加起来约占多少？（别忘了四种碱基加起来是 100\%。）
\item 一段 20 个碱基的 DNA 序列，最多能有多少种不同写法？这个数大约是全世界人口（约 $8\times 10^{9}$）的多少倍？
\item 关于照片 51 号：如果你是威尔金斯，正确的做法是什么？和同学讨论：“竞争”和“遵守规则”矛盾吗？再查一查富兰克林的生平，看看今天的人们怎样纪念她。
\end{sikaoti}
